\lyseis{Α3.5}{Διαίρεση ρητών}
\sol{1} α) $-3$ \quad β) $-4$ \quad γ) $-9$ \quad δ) $+70$.
\sol{2} α) $\frac{2}{5}\cdot\frac{15}{4}=\frac{30}{20}=\frac{3}{2}$ \quad β) $\frac{1}{2}\cdot(-3)=-\frac{3}{2}$ \quad γ) $4\cdot2=8$ \quad δ) $\frac{4}{3}\cdot\frac{1}{8}=\frac{4}{24}=\frac{1}{6}$ \quad ε) $\frac{18}{5}\cdot\frac{5}{6}=3$ \quad στ) $\frac{12}{21}\cdot\frac{4}{7}=\frac{48}{147}=\frac{16}{49}$.
\sol{3} α) $\frac{7}{11}\cdot\frac{22}{3}=\frac{14}{3}$ \quad β) $\frac{27}{49}\cdot\frac{14}{9}=\frac{6}{7}$ \quad γ) $4\cdot\frac{9}{56}=\frac{9}{14}$ \quad δ) $\frac{11}{2}\cdot\frac{1}{121}=\frac{1}{22}$ \quad ε) $2\frac{1}{8}=\frac{17}{8}$, άρα $\frac{3}{4}\cdot\frac{8}{17}=\frac{6}{17}$.
\sol{4} α) $450:15=30$ \quad β) $0{,}58$ (γιατί $49\cdot0{,}58=28{,}42$) \quad γ) $357{,}6:15=23{,}84$ \quad δ) $10$ (πολλαπλασιάζω επί $1\,000$: $292\,620:29\,262=10$) \quad ε) $4{,}8:6=0{,}8$.
\sol{5} α) $4{,}57$ \quad β) $0{,}032$ \quad γ) $0{,}0005$ \quad δ) $64$ \quad ε) $35$ \quad στ) $12\,000$.
\sol{6} α) Σωστό. β) Λάθος. γ) Σωστό. δ) Λάθος: $(-12):(-4)=+3$. ε) Λάθος: $3:\frac{1}{3}=3\cdot3=9$. στ) Λάθος: π.χ. $3:\frac{1}{2}=6>3$.
\sol{7} α) $-\frac{3}{4}+\frac{5}{8}+\frac{1}{2}=-\frac{6}{8}+\frac{5}{8}+\frac{4}{8}=\frac{3}{8}$. β) $1-\frac{5}{4}+\frac{1}{8}=\frac{8}{8}-\frac{10}{8}+\frac{1}{8}=-\frac{1}{8}$.
\sol{8} α) $x=12:(-3)=-4$ \quad β) $x=4:\frac{2}{3}=4\cdot\frac{3}{2}=6$ \quad γ) $x=\left(-\frac{1}{2}\right):(-5)=\frac{1}{10}$.
\sol{9} $6:\frac{3}{4}=6\cdot\frac{4}{3}=8$ μπουκάλια.
\sol{10} $4{,}5:0{,}15=450:15=30$ κομμάτια.
\sol{11} $6:0{,}75=600:75=8$. Το κιλό κοστίζει $8$ €.
\sol{12} $(-24{,}60):6=-4{,}10$. Κάθε παιδί έχει μερίδιο $-4{,}10$ €, δηλαδή χρωστάει $4{,}10$ €.
\sol{13} $\frac{3}{5}:\frac{1}{10}=\frac{3}{5}\cdot10=6$. Τρέχει $6$ χιλιόμετρα σε μία ώρα.
\sol{14} Π.χ. $3:\frac{1}{2}=6>3$. Το πηλίκο είναι μεγαλύτερο από τον (θετικό) διαιρετέο όταν ο διαιρέτης είναι θετικός και μικρότερος από το $1$, γιατί τότε ο διαιρέτης «χωράει» περισσότερες φορές.
\sol{15} Αν υπήρχε $x=5:0$, θα έπρεπε $x\cdot0=5$. Όμως κάθε αριθμός επί $0$ δίνει $0$, άρα δεν υπάρχει τέτοιος αριθμός.
\sol{16} $3:\frac{1}{2}=6$ και $3:\frac{1}{3}=9$. Το $\frac{1}{3}$ χωράει περισσότερες φορές στο $3$ (τρία τρίτα σε κάθε ένα), άρα το $3:\frac{1}{3}$ είναι μεγαλύτερο.
\sol{17} Αν $\frac{\alpha}{\beta}>0$, οι $\alpha,\beta$ είναι ομόσημοι (και οι δύο θετικοί ή και οι δύο αρνητικοί). Αν $\frac{\alpha}{\beta}<0$, είναι ετερόσημοι.
\sol{18} Αναμενόμενο: $1{,}2:0{,}15=120:15=8$ κομμάτια. Ο μαθητής ελέγχει κόβοντας.
\sol{19} Ελεύθερη απάντηση. Π.χ. $\dfrac{\frac{1}{2}}{\frac{3}{4}}=\frac{1}{2}\cdot\frac{4}{3}=\frac{2}{3}$.
